\documentclass[../../../include/open-logic-section]{subfiles}

\begin{document}

\olfileid{sth}{ord-arithmetic}{using-addition}
\olsection{کاربردِ جمع ترتیبی}

با استفاده از جمع اعداد ترتیبی، می‌توانیم رتبهٔ مجموعه‌های گوناگون
را به معنای \olref[spine][rank]{defnsetrank} به‌صراحت محاسبه کنیم:

\begin{lem}\ollabel{rankcomputation}
اگر $\setrank{A} = \alpha$ و $\setrank{B} = \beta$، آن‌گاه:
\begin{enumerate}
	\item\ollabel{exrankpow} $\setrank{\Pow{A}} = \alpha\ordplus 1$
	\item\ollabel{exrankpair} $\setrank{\{A, B\}} = \max(\alpha,
	\beta) \ordplus 1$
	\item\ollabel{exrankcup} $\setrank{A \cup B} = \max(\alpha,
	\beta)$
	\item\ollabel{exranktuple} $\setrank{\tuple{A,B}} = \max(\alpha,
	\beta) \ordplus  2$
	\item\ollabel{exranktimes} $\setrank{A \times B} \leq \max(\alpha,
	\beta) \ordplus  2$
	\item\ollabel{exrankunion} اگر $\alpha$ تهی یا حدی باشد،
	$\setrank{\bigcup A} = \alpha$؛ و اگر
	$\alpha = \gamma\ordplus 1$، آن‌گاه
	$\setrank{\bigcup A} = \gamma$
\end{enumerate}
\end{lem}

\begin{proof}
در سراسر اثبات بارها
\olref[spine][rank]{ranksupstrict} را به کار می‌بریم.

\emph{\olref{exrankpow}.} اگر $x \subseteq A$، آن‌گاه
$\setrank{x} \leq \setrank{A}$. پس
$\setrank{\Pow{A}} \leq \alpha \ordplus 1$. از آنجا که به‌ویژه
$A \in \Pow{A}$، داریم
$\setrank{\Pow{A}} = \alpha \ordplus 1$.

\emph{\olref{exrankpair}.} بنا بر
\olref[spine][rank]{ranksupstrict}.

\emph{\olref{exrankcup}.} بنا بر
\olref[spine][rank]{ranksupstrict}.

\emph{\olref{exranktuple}.} با دو بار استفاده از
\olref{exrankpair}.

\emph{\olref{exranktimes}.} توجه کنید که
$A \times B \subseteq \Pow{\Pow{A \cup B}}$؛ سپس
\olref{exrankpow} را دو بار و \olref{exrankcup} را یک بار به کار
ببرید.

\emph{\olref{exrankunion}.} اگر $\alpha = \gamma\ordplus 1$، عضوی
مانند $c \in A$ هست که $\setrank{c} = \gamma$، و هیچ
!!{element} از $A$ رتبهٔ بالاتری ندارد؛ پس
$\setrank{\bigcup A} = \gamma$. اگر $\alpha$ عدد ترتیبیِ حدی باشد،
$A$ دارای !!{element}هایی با رتبه‌های دلخواه نزدیک به $\alpha$
(اما اکیداً کوچک‌تر از آن) است؛ پس $\bigcup A$ نیز دارای
!!{element}هایی با رتبه‌های دلخواه نزدیک به $\alpha$ (اما اکیداً
کوچک‌تر از آن) است، و بنابراین
$\setrank{\bigcup A} = \alpha$.
\end{proof}
\noindent
این‌که چرا \olref{exranktimes} یک \emph{نامساوی} دارد، به‌عنوان
تمرین واگذار می‌شود.

\begin{prob}
مجموعه‌های $A$ و $B$ای بیابید که
$\setrank{A \times B} =
\max(\setrank{A}, \setrank{B})$. مجموعه‌های $A$ و $B$ای بیابید که
$\setrank{A \times B} =
\max(\setrank{A}, \setrank{B}) \ordplus 2$. آیا رتبه‌های دیگری نیز
ممکن‌اند؟
\end{prob}

اکنون همچنین می‌توانیم نشان دهیم چند معنای معقول برای «متناهی» یا
«نامتناهی» خواندن یک عدد ترتیبی با یکدیگر منطبق‌اند:

\begin{lem}\ollabel{ordinfinitycharacter}
برای هر عدد ترتیبیِ $\alpha$، گزاره‌های زیر هم‌ارزند:
\begin{enumerate}
	\item\ollabel{ord:notinomega} $\alpha\notin \omega$؛ یعنی
	$\alpha$ عدد طبیعی نیست؛
	\item\ollabel{ord:omegaplus} $\omega \leq \alpha$؛
	\item\ollabel{ord:oneplus} $1 \ordplus \alpha = \alpha$؛
	%\alpha \approx \alpha \ordplus 1$
	\item\ollabel{ord:plusone} $\alpha \approx \alpha\ordplus 1$؛
	یعنی $\alpha$ و $\alpha\ordplus 1$ هم‌شمارند؛
	\item\ollabel{ord:infinite} $\alpha$ نامتناهیِ ددکیندی است.
	\end{enumerate}
\end{lem}
\noindent
پس پنج شیوهٔ اثباتاً هم‌ارز برای فهم متناهی یا نامتناهی‌بودن یک عدد
ترتیبی داریم.

\begin{proof}
\emph{\olref{ord:notinomega} $\Rightarrow$ \olref{ord:omegaplus}.}
بنا بر سه‌شقی.

\emph{\olref{ord:omegaplus} $\Rightarrow$ \olref{ord:oneplus}.}
$\alpha \geq \omega$ را ثابت کنید. بنا بر استقرای ترامتناهی، کوچک‌ترین
عدد ترتیبیِ $\gamma$ (شاید $0$) وجود دارد به‌طوری‌که برای یک عدد
ترتیبیِ حدیِ $\beta$ داشته باشیم
$\alpha = \beta \ordplus \gamma$. اکنون:
\[
	1 \ordplus \alpha =
	1 \ordplus (\beta \ordplus \gamma) =
	(1 \ordplus \beta) \ordplus \gamma =
	\supstrict_{\delta < \beta} (1 \ordplus  \delta) \ordplus  \gamma =
	\beta \ordplus  \gamma =
	\alpha.
\]
\emph{\olref{ord:oneplus} $\Rightarrow$ \olref{ord:plusone}.}
آشکارا !!a{bijection} مانند
$f \colon (\alpha \disjointsum 1) \to (1 \disjointsum \alpha)$ وجود
دارد. اگر $1 \ordplus \alpha = \alpha$، یکریختی مانند
$g \colon (1 \disjointsum \alpha) \to \alpha$ وجود دارد. اکنون
$\comp{f}{g}$ را در نظر بگیرید.

\emph{\olref{ord:plusone} $\Rightarrow$ \olref{ord:infinite}.}
اگر $\alpha \approx \alpha \ordplus 1$، !!a{bijection} مانند
$f \colon (\alpha \disjointsum 1) \to \alpha$ وجود دارد. برای هر
$\gamma < \alpha$ تعریف کنید
$g(\gamma) = f(\gamma, 0)$؛ این !!{injection} گواه نامتناهیِ
ددکیندی‌بودن $\alpha$ است، زیرا
$f(0,1) \in \alpha \setminus \ran{g}$.

\emph{\olref{ord:infinite} $\Rightarrow$ \olref{ord:notinomega}.}
این همان \olref[z][infinity-again]{naturalnumbersarentinfinite} است.
\end{proof}

\end{document}
